Write $s_{k} = \sqrt{m_{k}}$, $S_{1} = \sum_{k}s_{k}$, $S_{2} = \sum_{k}s_{k}^{2}$, so $Q = S_{2}/S_{1}^{2}$, the quotient of Proposition \claimref{families.linear-conjecture-refuted-data}.

\emph{(i)} $m \mapsto \lambda m$ sends $S_{2} \mapsto \lambda S_{2}$ and $S_{1} \mapsto \sqrt{\lambda}S_{1}$, so $Q$ is unchanged. The check confirms it exactly on rational square coordinates and confirms that a factor on one mass alone changes $Q$.

\emph{(ii)} $\partial Q/\partial s_{j} = (2s_{j}S_{1}^{2} - 2S_{1}S_{2})/S_{1}^{4} = 2(S_{1}s_{j} - S_{2})/S_{1}^{3}$; the check compares it with exact difference quotients. By Cauchy-Schwarz $S_{1}^{2} \le 3S_{2}$, and $S_{1} > s_{3}$, so $S_{2} \ge S_{1}^{2}/3 > S_{1}s_{3}/3$; hence $3s_{j} \le s_{3}$ gives $S_{1}s_{j} < S_{2}$ and $\partial Q/\partial s_{j} < 0$. For $j = 3$, $S_{2} = \sum s_{k}^{2} < s_{3}S_{1}$ since $s_{k} < s_{3}$ for $k \ne 3$ (or $\le$ with a strict term), so $\partial Q/\partial s_{3} > 0$. By (i) take $D_{3} = 1$. Along $t \mapsto (D_{1}^{t}m_{1}, D_{2}^{t}m_{2}, m_{3})$, $t \in [0,1]$, each light root rises monotonically to $\sqrt{D_{j}m_{j}} \le s_{3}/3$, so the condition $3s_{j}(t) \le s_{3}$ holds throughout and $dQ/dt = \sum_{j=1,2}(\tfrac12\ln D_{j})\,s_{j}(t)\,\partial Q/\partial s_{j} \le 0$, strictly unless $D_{1} = D_{2} = 1$. So $Q(m') \le Q(m)$. The check also confirms it on a rational grid of factors.

\emph{(iii)} The values (GeV): down stored $m_{d} = 0.00470$, $m_{s} = 0.0935$ (at $2$ GeV), $m_{b} = 4.18$ ($m_{b}(m_{b})$), from the stored quantities; up $m_{t} = 172.57$ (stored, direct measurements). Cited from memory and flagged: $m_{u} = 0.00216$ (at $2$ GeV), $m_{c}(m_{c}) = 1.273$, the charm pole mass $1.67$, and the running masses at $M_{Z}$, $(m_{u}, m_{c}, m_{t}) = (0.00127, 0.619, 171.7)$ and $(m_{d}, m_{s}, m_{b}) = (0.00290, 0.055, 2.89)$ (the 2008 table of Xing, Zhang and Zhou, Phys. Rev. D 77 (2008) 113016, recalled, not checked against the source). With certified rational bounds on every square root (integer square roots at $10^{-12}$), $3\sum m > 2(\sum\sqrt{m})^{2}$ holds in all five sets, with $Q \in (0.731, 0.732)$, $(0.745, 0.746)$, $(0.848, 0.849)$, $(0.831, 0.832)$, $(0.888, 0.889)$. If $m = D m^{0}$ with $D$ of kind (i), $Q(m^{0}) = Q(m)$; of kind (ii), $Q(m^{0}) \ge Q(m)$ by (ii) applied to $m^{0}$ (the hypothesis $D_{j}m^{0}_{j} \le D_{3}m^{0}_{3}/9$ is $m_{j} \le m_{3}/9$, true in every set). In both cases $Q(m^{0}) > \tfrac23$.

\emph{(iv)} $m^{0}_{k} = m_{k}/D_{k}$; by (i) rescale so that $m^{0} = (\rho_{1}m_{1}, \rho_{2}m_{2}, m_{3})$ with $\rho_{j} = D_{3}/D_{j} \le R$. By (ii) (its derivative statement, applied with the light factors moved up to $R$; the condition $Rm_{j} \le m_{3}/9$ is checked at the upper ends used) $Q(m^{0}) \ge Q(Rm_{1}, Rm_{2}, m_{3})$, whatever the $\rho_{j} \in (0, R]$. The check certifies $Q(R_{\mathrm{lo}}m_{1}, R_{\mathrm{lo}}m_{2}, m_{3}) > \tfrac23 > Q(R_{\mathrm{hi}}m_{1}, R_{\mathrm{hi}}m_{2}, m_{3})$ for the four stated brackets; $Q$ at the corner falls strictly with $R$, so $R^{*}$, its crossing of $\tfrac23$, lies between them (computed: $1.889$, $2.197$, $8.675$, $17.56$: $\sqrt{R^{*}}$ is the smaller root of $R(3p - 2a^{2}) - 4ab\sqrt{R} + m_{3} = 0$ with $a = s_{1} + s_{2}$, $b = s_{3}$, $p = m_{1} + m_{2}$).

\emph{The strong dressing (argued; standard QCD cited from memory, not stored).} In the MS-bar scheme the quark mass anomalous dimension depends on $\alpha_{s}$ and the number of active flavours but not on the flavour, and the decoupling at a heavy threshold rescales all lighter masses alike, so the running between two common scales is a factor of kind (i); this is why $Q$ at $M_{Z}$ is the bracket's $Q$ for any bracket defined at a common scale. Quoting each mass at its own scale $\mu_{k}$ increasing with $m_{k}$ multiplies it by $m(\mu_{k})/m(\mu_{0})$, which decreases with $\mu_{k}$ (asymptotic freedom), and the pole conversion $m_{\mathrm{pole}} = \bar m(\bar m)(1 + \tfrac{4}{3\pi}\alpha_{s}(\bar m) + \dots)$ has a positive correction decreasing with the mass: both favour the lighter quark, kind (ii) (the masses involved satisfy the ninth condition). Constituent dressing, of order $\Lambda_{\mathrm{QCD}}$ additive, favours the lighter quarks still more, and the top decays before it hadronises. So no strong dressing of these kinds can raise a Koide bracket to the measured quotients, and the selectivity account of Observation \claimref{families.selectivity-observation-level}, as far as it attributes the quark failure to the strong dressing, is excluded unless a dressing favouring the heaviest quark by at least $R^{*}$ is exhibited.
